Instruction tuning rewrites a base model's next-token distribution, but most of that change is anticipatable. We ask a sharper question than prior work: where does alignment change the distribution in ways that a learned predictor cannot anticipate from simple base-side context, and what kinds of tokens live there? The Superficial Alignment Hypothesis holds that alignment mostly adjusts surface style, but no one has trained a predictor of per-token divergence and studied the errors it makes. We do exactly that. For each instruction we greedily decode an instruct response, teacher-force the identical sequence through both base and instruct models under the same chat-formatted context to isolate the weight difference, and at every response position compute the target divergence plus cheap base-side features (position, base entropy, base top-1 probability, emitted-token log-probability, token class). We fit a gradient-boosted predictor on one corpus and read its residuals on a disjoint one, across three Qwen2.5 scales (0.5B, 1.5B, 7B). The predictor explains roughly 60--68\% of per-token divergence cross-corpus, driven almost entirely by base-side uncertainty; position alone is nearly useless. The residual, where the model is provably wrong, is not noise: it concentrates significantly on safety prompts and on discourse and structural pivots, persona and identity tokens (the model naming its maker, framing a knowledge cutoff), and refusal-recovery moves. This refines the Superficial Alignment Hypothesis: the average shift is predictable and largely superficial, but the unpredictable residual is substantive, alignment-specific, and localizable with no labels.
